\chapter{Harvesting the Variance Risk Premium}\label{ch:s2:harvesting-the-variance-risk-premium}

At the end of 84\% of months since 1990 the VIX index, the market's price of the next month's S\&P 500 volatility, was higher than the volatility that
followed: on average 19.5 against 15.4. Selling that price and delivering the realised volatility has been one of the most reliable trades in finance.
Its losses come the other way, all at once. At the end of February 2020 the VIX stood at 40.1; the next month's realised volatility was 94.5. On this
chapter's synthetic market a short-variance book earns a Sharpe ratio of 1.57 until its worst month, which then costs 7.7 of its monthly standard
deviations: 22\% of capital at a 10\% volatility target, and more than the capital at 60\%. This chapter is about the implementations of short
volatility and about sizing them for the month that matters. The build is two components: \texttt{firm.synthvol}, the synthetic option market that
Part~I runs on, and \texttt{firm.shortvol}.

\section{The premium and why it exists}

The variance risk premium (Book~5, chapter~14) is the difference between the price of variance, the rate of a variance swap or the square of an index
like the VIX, and the variance that is then realised. Carr and Wu measured it directly: they synthesised variance swap rates from option prices on
five stock indexes and 35 stocks, and took the difference with realised variance as the premium. It is paid because variance is highest when
investors are poorest: a long-variance position pays in crashes, and investors pay to hold it, as they pay for insurance.

\begin{definition}[Short-volatility strategy]\label{def:s2:harvesting-the-variance-risk-premium:short}
A \emph{short-volatility strategy}\index{short-volatility strategy} earns the variance risk premium by selling options or variance and delivering
realised volatility: short variance swaps, delta-hedged short options, covered calls, short puts, or short volatility futures. Its returns are small and
frequent in calm markets and large and negative when volatility jumps.
\end{definition}

The real record, from Cboe's data (\cref{fig:s2:harvesting-the-variance-risk-premium:gap}): over 440 month-ends from January 1990 to August 2026 the VIX
averaged 19.5\% and the next month's realised volatility of the S\&P 500 15.4\%; the VIX was higher 84\% of the time, by 4.1 points on average; its
square exceeded realised variance by 31\% on average. The six months when it fell furthest short tell the other side: February 2020 (VIX 40.1, realised
94.5), September 2008 (39.4 and 82.3), August 2008 (20.6 and 54.3), March 2025 (22.3 and 47.7), July 2011 and June 2002.

\begin{omfigure}
\begin{tikzpicture}
\begin{axis}[width=11cm,height=5.2cm,xlabel={\small year},ylabel={\small VIX minus realised (points)},tick label style={font=\footnotesize},
  xmin=1990,xmax=2027,ymin=-15,ymax=15,grid=major,grid style={gray!25},table/col sep=comma,/pgf/number format/1000 sep={}]
\addplot[omThm,thick] table[x=year,y=gap_pts]{figdata/strategies-2/01-harvesting-the-variance-risk-premium/gap.csv};
\addplot[black!40,dashed] coordinates {(1990,0) (2027,0)};
\end{axis}
\end{tikzpicture}

\omcaption{The S\&P 500 variance premium, 1990--2026: the twelve-month average, at each month-end, of the VIX minus the realised volatility of the S\&P
500 over the following 21 trading days. Derived from Cboe's VIX and S\&P 500 index histories; the raw series are not redistributed. Data:
\texttt{s2\_fetch\_cboe}.}\label{fig:s2:harvesting-the-variance-risk-premium:gap}
\end{omfigure}

\section{Implementations: variance, options, overwriting}

\texttt{firm.synthvol} (\cref{lst:s2:harvesting-the-variance-risk-premium:synthvol}) simulates twenty years of an index and thirty member stocks: a common
factor with stochastic variance around 16\% (mean reversion of four a year, a leverage correlation of $-0.7$), negative jumps once a year on average, and
three planted crashes of 25\%, 20\% and 30\% over ten days each. The truth sets the surface: the index's at-the-money implied variance is the expected
variance over the option's life times 1.3, and the smile falls with strike. Realised volatility over the twenty years is 19.0\%, crashes included, and
the average one-month implied volatility 18.3\%.

\begin{definition}[Option overwriting]\label{def:s2:harvesting-the-variance-risk-premium:overwrite}
\emph{Option overwriting}\index{option overwriting} sells options against a position already held, most often calls against a long equity position
(covered calls), collecting premium in exchange for giving up gains beyond the strike; it is a \omterm{def:s2:harvesting-the-variance-risk-premium:short}{short-volatility strategy} with the market's direction
added.
\end{definition}

\texttt{firm.shortvol} runs five monthly implementations on the index, each rolled every 21 trading days:

\begin{center}
\small
\begin{tabular}{@{}lcccc@{}}
\toprule
implementation, 239 months & Sharpe ratio & skewness & worst month (sd) & months up\\
\midrule
short variance swap & 0.69 & $-3.53$ & $-7.7$ & 79\%\\
delta-hedged short straddle & 1.24 & $-2.32$ & $-6.4$ & 77\%\\
covered calls, 2\% out of the money & 0.40 & $-4.24$ & $-8.1$ & 69\%\\
short puts, 5\% out of the money & 0.08 & $-7.22$ & $-10.2$ & 92\%\\
short variance sized by $1/\text{strike}$ & 0.09 & $-12.45$ & $-14.2$ & 79\%\\
the index itself & 0.26 & $-1.82$ & $-6.8$ & 62\%\\
\bottomrule
\end{tabular}
\end{center}

Every implementation has negative skewness, and the shape matches the premium's source: the short put wins in 92\% of months and loses ten standard
deviations in one. Sizing short variance by the inverse of its strike (more when volatility is cheap) makes it worse: the planted crashes start from
calm markets, when the book is largest. The synthetic crashes are harsher than recent history: on Cboe's indices from February 2007 to September 2026,
PutWrite returned 7.2\% a year at 13.8\% volatility (a Sharpe ratio of 0.57 without a risk-free rate) and BuyWrite 6.0\% at 14.3\% (0.48), against
9.0\% at 19.7\% (0.53) for the S\&P 500 price index, which omits dividends. PutWrite lost 27.6\% from September to November 2008 and 19.8\% in February
and March 2020; BuyWrite 28.5\% and 21.3\%.

\section{Delta-hedged option returns}

\begin{definition}[Delta-hedged option return]\label{def:s2:harvesting-the-variance-risk-premium:hedged}
A \emph{delta-hedged option return}\index{delta-hedged option return} is the gain of an option position whose delta is hedged in the underlying at
regular intervals, relative to its premium; under a model with no volatility risk premium its expectation is close to zero, so its sign and size measure
the premium.
\end{definition}

Bakshi and Kapadia used exactly this: in S\&P 500 options, a delta-hedged long option underperformed zero, less so away from the money, more when
volatility was high, even after accounting for jump fears. The seller's side, a short at-the-money straddle hedged daily at the surface's current
volatility (\cref{lst:s2:harvesting-the-variance-risk-premium:shortvol}), earns 0.53\% of the index a month on the synthetic market, about an eighth of the
straddle's premium of 4.2\%, with a Sharpe ratio of 1.24, the best of the five, because the hedge removes the index's own direction and leaves the gap between
implied and realised volatility.

\section{Blow-ups and how to size for them}

Short volatility is a sizing problem. \Cref{fig:s2:harvesting-the-variance-risk-premium:paths} scales the short-variance book to three volatility
targets. Its Sharpe ratio before the worst month (the first crash, in year 5.9) is 1.57 at any size; the worst month costs 7.7 of its monthly standard
deviations at any size:

\begin{center}
\small
\begin{tabular}{@{}lccccc@{}}
\toprule
target volatility & 5\% & 10\% & 20\% & 40\% & 60\%\\
\midrule
Sharpe ratio before the worst month & 1.57 & 1.57 & 1.57 & 1.57 & 1.57\\
worst month, share of capital & $-11.1\%$ & $-22.2\%$ & $-44.3\%$ & $-88.6\%$ & $-132.9\%$\\
capital survives & yes & yes & yes & yes & no\\
\bottomrule
\end{tabular}
\end{center}

The record before the crash says nothing about the size of the crash: a manager who judged the risk from five years of monthly returns saw a Sharpe
ratio above 1.5 and a worst month of a fraction of the one to come. The size that survives is set by the stress, not the volatility: a book that must
not lose more than a quarter of its capital in a month can run the synthetic short-variance book at about 10\% volatility, whatever its history.

\begin{omfigure}
\begin{tikzpicture}
\begin{axis}[width=11cm,height=5.4cm,ymode=log,xlabel={\small year},ylabel={\small capital (log scale)},tick label style={font=\footnotesize},
  xmin=0,xmax=20,ymin=0.3,ymax=50,ytick={0.5,1,2,5,10,20,50},yticklabels={0.5,1,2,5,10,20,50},grid=major,grid style={gray!25},legend columns=3,legend style={font=\scriptsize,at={(0.5,-0.3)},anchor=north,
  draw=none,fill=none,/tikz/every even column/.append style={column sep=8pt}},table/col sep=comma]
\addplot[omExo,thick] table[x=year,y=v10]{figdata/strategies-2/01-harvesting-the-variance-risk-premium/paths.csv};
\addplot[omDef,thick] table[x=year,y=v20]{figdata/strategies-2/01-harvesting-the-variance-risk-premium/paths.csv};
\addplot[omThm,thick] table[x=year,y=v40]{figdata/strategies-2/01-harvesting-the-variance-risk-premium/paths.csv};
\legend{10\% target,20\% target,40\% target}
\end{axis}
\end{tikzpicture}

\omcaption{The synthetic short-variance book scaled to three volatility targets, compounded monthly over twenty years; the planted crashes of years 6.0 and
17.5 show as single-month falls, the smaller one of year 12.7 barely. Data: \texttt{s2\_vrp.paths}.}\label{fig:s2:harvesting-the-variance-risk-premium:paths}
\end{omfigure}

\section{Strategy files}

\begin{strategyfile}{Short index variance}\label{strat:s2:harvesting-the-variance-risk-premium:variance}
\sfield{Who pays you, and why} Investors who pay for protection against variance, which is highest in bad times.
\sfield{Instruments and venues} Variance swaps over the counter; a strip of listed options.
\sfield{Signal} The variance premium: implied variance against a forecast of realised variance.
\sfield{Sizing and execution} Sized by a stress (the worst plausible month), not by volatility; monthly rolls.
\sfield{Costs} Option spreads on the replicating strip; dealer margins on swaps.
\sfield{How it dies} Crashes; crowding, which raises the crash's size.
\sfield{Horizon, capacity, infrastructure} A month; large capacity in index variance; a surface and stress system.
\sfield{Backtest honestly} Crash years in the sample; variance strikes as quoted; margin calls.
\sfield{Sources} Carr and Wu (2009); Cboe VIX data; this chapter: 1.57 before the worst month, $-7.7$ standard deviations in it.
\end{strategyfile}

\begin{strategyfile}{Delta-hedged short straddle}\label{strat:s2:harvesting-the-variance-risk-premium:straddle}
\sfield{Who pays you, and why} As for short variance, with the direction hedged away.
\sfield{Instruments and venues} Listed at-the-money index options and futures for the hedge.
\sfield{Signal} Implied against forecast volatility.
\sfield{Sizing and execution} Sold monthly; delta hedged daily or on bands.
\sfield{Costs} Option spreads and hedge trading.
\sfield{How it dies} Gap moves the hedge cannot follow.
\sfield{Horizon, capacity, infrastructure} A month; a hedging engine.
\sfield{Backtest honestly} Hedges at executable prices; gaps and halts.
\sfield{Sources} Bakshi and Kapadia (2003); this chapter: a Sharpe ratio of 1.24 on the synthetic index.
\end{strategyfile}

\begin{strategyfile}{Covered-call overwriting}\label{strat:s2:harvesting-the-variance-risk-premium:overwrite}
\sfield{Who pays you, and why} Buyers of upside calls; the variance premium on calls.
\sfield{Instruments and venues} An equity holding and listed calls on it.
\sfield{Signal} None beyond the premium; strike and tenor chosen by rule.
\sfield{Sizing and execution} Calls written against the whole holding each month.
\sfield{Costs} Option spreads; the upside given up.
\sfield{How it dies} It does not die; it lags in strong rallies and falls with the market.
\sfield{Horizon, capacity, infrastructure} A month; very large capacity.
\sfield{Backtest honestly} Include dividends in the comparison index.
\sfield{Sources} Cboe BuyWrite index: 6.0\% a year at 14.3\% volatility, 2007--2026.
\end{strategyfile}

\begin{strategyfile}{Put writing}\label{strat:s2:harvesting-the-variance-risk-premium:putwrite}
\sfield{Who pays you, and why} Buyers of crash protection, who pay most for out-of-the-money puts.
\sfield{Instruments and venues} Listed index puts, collateralised with cash.
\sfield{Signal} None beyond the premium.
\sfield{Sizing and execution} Fully collateralised; one-month puts rolled at expiry.
\sfield{Costs} Option spreads.
\sfield{How it dies} Crashes: the short put wins most months and loses in a few.
\sfield{Horizon, capacity, infrastructure} A month; very large capacity.
\sfield{Backtest honestly} Collateral earning the bill rate; crash years in the sample.
\sfield{Sources} Cboe PutWrite index: 7.2\% a year at 13.8\% volatility, 2007--2026; this chapter's harsher synthetic crashes.
\end{strategyfile}

\begin{strategyfile}{Volatility-targeted short volatility}\label{strat:s2:harvesting-the-variance-risk-premium:voltarget}
\sfield{Who pays you, and why} As for short variance.
\sfield{Instruments and venues} Variance swaps or straddles.
\sfield{Signal} The premium, sized by current implied volatility.
\sfield{Sizing and execution} Larger when implied volatility is low, smaller when high.
\sfield{Costs} Extra turnover.
\sfield{How it dies} Crashes that start from calm markets, when the book is largest.
\sfield{Horizon, capacity, infrastructure} A month.
\sfield{Backtest honestly} Crashes from low-volatility starts in the sample.
\sfield{Sources} This chapter: sizing by the inverse of the strike cut the Sharpe ratio from 0.69 to 0.09.
\end{strategyfile}

\section{Tutorial: picking up nickels}\label{tut:s2:harvesting-the-variance-risk-premium:tut}

\begin{tutorial}
\textbf{Goal.} Build the synthetic option market, run five short-volatility implementations, and size the short-variance book through its planted
crashes; compare with the real variance premium and option-writing indices. \textbf{End state:} the tables and the two figures.
\begin{tutsteps}
\item \textbf{The market}: stochastic variance with leverage, jumps and crashes; members with specific returns.
\omcode{code/firm/synthvol/firm_synthvol.py}{95}{120}{Twenty years of the synthetic index and its members.}\label{lst:s2:harvesting-the-variance-risk-premium:synthvol}
\item \textbf{Two implementations}: the variance swap and the delta-hedged straddle.
\omcode{code/firm/shortvol/firm_shortvol.py}{33}{58}{Short variance and a delta-hedged short straddle.}\label{lst:s2:harvesting-the-variance-risk-premium:shortvol}
\item \textbf{Run} \texttt{s2\_fetch\_cboe.py} once for the real statistics, then \texttt{implementations()}, \texttt{leverage\_table()} and
\texttt{fig\_vrp.py}.
\end{tutsteps}
\textbf{What to change next.} Hedge the straddle on bands instead of daily; add a crash that starts from high volatility and see whether inverse-strike
sizing helps; buy a far out-of-the-money put against the short variance and measure what it costs and saves.
\end{tutorial}

\section{Build: the synthetic option market and short volatility}\label{bld:s2:harvesting-the-variance-risk-premium:shortvol}

\begin{build}
\bfield{Purpose} \texttt{firm.synthvol}: a deterministic index and member option market with stochastic variance, jumps, crashes and surfaces carrying
a variance premium and a correlation premium, for Part~I. \texttt{firm.shortvol}: short-volatility implementations and leverage.
\bfield{Interface} \texttt{VolConfig(\dots)}, \texttt{simulate\_vol(cfg)}, \texttt{expected\_var}, \texttt{atm\_iv}, \texttt{smile\_iv}, \texttt{bs\_price},
\texttt{bs\_delta}; \texttt{periods}, \texttt{short\_variance}, \texttt{hedged\_straddle}, \texttt{overwrite}, \texttt{put\_write}, \texttt{lever}.
\bfield{Rules} Implied volatilities from the truth known at the time; options priced at the surface; P\&L per unit of index at inception.
\bfield{Acceptance tests} \texttt{code/firm/synthvol/tests/} and \texttt{code/firm/shortvol/tests/}: Black--Scholes by hand, expected variance and the
premium, a planted crash; variance and straddle P\&L on a flat path, a wipe-out.
\bfield{Stretch} Surfaces with term-structure premia; transaction costs on hedges; members' options.
\end{build}

\begin{omsources}
\item P.~Carr and L.~Wu, ``Variance risk premiums'', \emph{Review of Financial Studies} 22(3), 2009.
\item G.~Bakshi and N.~Kapadia, ``Delta-hedged gains and the negative market volatility risk premium'', \emph{Review of Financial Studies} 16(2), 2003.
\item Cboe Global Markets, daily histories of the VIX, S\&P 500, PutWrite and BuyWrite indices.
\end{omsources}

\section{Exercises}

\begin{exercise}[$\star$]\label{exo:s2:harvesting-the-variance-risk-premium:1}
An at-the-money straddle is worth about $0.8\,\sigma\sqrt{\tau}$ of the index. What is it worth for a month (21 trading days) at 20\% volatility?
\end{exercise}

\begin{exercise}[$\star$]\label{exo:s2:harvesting-the-variance-risk-premium:2}
Expected variance is $0.16^2$ and implied variance is 30\% higher. What are the implied variance and volatility?
\end{exercise}

\begin{exercise}[$\star$]\label{exo:s2:harvesting-the-variance-risk-premium:3}
Why does the short put win more often than the short variance swap but have the lower Sharpe ratio?
\end{exercise}

\begin{exercise}[$\star\star$]\label{exo:s2:harvesting-the-variance-risk-premium:4}
A book at 10\% annual volatility loses 7.7 monthly standard deviations. What share of capital does it lose?
\end{exercise}

\begin{exercise}[$\star\star$]\label{exo:s2:harvesting-the-variance-risk-premium:5}
Why does sizing short variance by the inverse of its strike make it worse on the synthetic market? When would it help?
\end{exercise}

\begin{exercise}[$\star\star$]\label{exo:s2:harvesting-the-variance-risk-premium:6}
Why does the VIX exceed realised volatility most of the time, and why is its average excess not the premium's full measure?
\end{exercise}

\begin{exercise}[$\star\star\star$]\label{exo:s2:harvesting-the-variance-risk-premium:7}
\emph{Coding.} Run \texttt{leverage\_table((0.3, 0.5))}. At which target volatility between them does the worst month take the whole capital?
\end{exercise}

\begin{exercise}[$\star\star\star$]\label{exo:s2:harvesting-the-variance-risk-premium:8}
\emph{Find the flaw.} ``Our short-volatility fund has a Sharpe ratio of 2 over five years and its worst month was $-3\%$; at 15\% volatility the risk
is modest.''
\end{exercise}

\section{Problem: Picking Up Nickels}

\begin{problem}[{Weekend problem --- short volatility, sized}]\label{pb:s2:harvesting-the-variance-risk-premium:1}
The chapter's synthetic option market, Cboe's data and the public record.

\medskip\noindent\textbf{Part I --- The premium.}
\begin{enumerate}
\item Define a \omterm{def:s2:harvesting-the-variance-risk-premium:short}{short-volatility strategy} and \omterm{def:s2:harvesting-the-variance-risk-premium:overwrite}{option overwriting}.
\item What did Carr and Wu measure, and how?
\item Give the VIX's record against realised volatility, 1990--2026.
\item Name the months when realised volatility exceeded the VIX by most.
\end{enumerate}

\medskip\noindent\textbf{Part II --- Implementations.}
\begin{enumerate}[resume]
\item Describe the synthetic market and how its surface is set.
\item Give the five implementations' Sharpe ratios and skewness.
\item Why is every skewness negative?
\item What did the PutWrite and BuyWrite indices earn, and how did they fare in 2008 and 2020?
\end{enumerate}

\medskip\noindent\textbf{Part III --- Hedged options.}
\begin{enumerate}[resume]
\item Define a \omterm{def:s2:harvesting-the-variance-risk-premium:hedged}{delta-hedged option return}.
\item What did Bakshi and Kapadia find?
\item What does the synthetic straddle earn, and why is its Sharpe ratio the highest?
\item What would a hedge on bands change?
\end{enumerate}

\medskip\noindent\textbf{Part IV --- The verdict.}
\begin{enumerate}[resume]
\item State the \emph{named result}: the short-variance book's Sharpe ratio before its worst month and its loss in that month, by leverage.
\item Why does the record before the crash not measure the crash?
\item How should a short-volatility book be sized?
\item Why did inverse-strike sizing fail here?
\item How would you backtest a \omterm{def:s2:harvesting-the-variance-risk-premium:short}{short-volatility strategy} honestly?
\item Which strategy file wins most often, and which has the best Sharpe ratio?
\item How does this chapter relate to Book~8, chapter~26's volatility targeting?
\item In one sentence: what does a short-volatility trader sell?
\end{enumerate}
\end{problem}

\section{Interview questions}

\begin{interviewq}[$\star$ \iqroles{researcher}]\label{iq:s2:harvesting-the-variance-risk-premium:1}
What is the variance risk premium, and why does it exist?
\end{interviewq}

\begin{interviewq}[$\star\star$ \iqroles{trader}]\label{iq:s2:harvesting-the-variance-risk-premium:2}
You are short a delta-hedged straddle and the index gaps down 5\% overnight. What happened to your P\&L, and what do you do?
\end{interviewq}

\begin{interviewq}[$\star\star$ \iqroles{risk}]\label{iq:s2:harvesting-the-variance-risk-premium:3}
How would you set limits on a short-volatility book?
\end{interviewq}

\begin{interviewq}[$\star\star$ \iqroles{researcher}]\label{iq:s2:harvesting-the-variance-risk-premium:4}
Compare selling variance swaps, straddles and out-of-the-money puts as ways to earn the premium.
\end{interviewq}

\begin{interviewq}[$\star\star$ \iqroles{developer}]\label{iq:s2:harvesting-the-variance-risk-premium:5}
What does a daily delta-hedging engine for an option book need?
\end{interviewq}

\begin{interviewq}[$\star\star\star$ \iqroles{researcher}]\label{iq:s2:harvesting-the-variance-risk-premium:6}
A delta-hedged option with gamma $\Gamma$ on an underlying at $S$ is hedged continuously. Show that its P\&L over a short interval is about
$\tfrac12\Gamma S^2(\sigma_r^2 - \sigma_i^2)\,dt$, with $\sigma_r$ realised and $\sigma_i$ implied volatility, and explain what this implies for its sign.
\end{interviewq}
